\documentclass[10pt]{letter}
\usepackage[a4paper,margin=22mm]{geometry}
\usepackage[T1]{fontenc}
\usepackage{lmodern}
\usepackage[hidelinks]{hyperref}
\usepackage{xcolor}
\setlength{\parindent}{0pt}
\setlength{\parskip}{3pt}
\pagestyle{empty}
% Marks numbers that depend on running experiments; resolve before submission.
\newcommand{\pending}[1]{\textcolor{red}{[PENDING: #1]}}
% Seed-1 value from the fair-baseline run; replace with the five-seed result.
\newcommand{\provisional}[1]{\textcolor{orange}{#1}}

\signature{Shishir Singh Chauhan\\Manipal University Jaipur}
\address{Shishir Singh Chauhan\\
Manipal University Jaipur\\
Dehmi Kalan, Jaipur-Ajmer Expressway\\
Jaipur 303007, Rajasthan, India}

\begin{document}

\begin{letter}{The Editors\\
Transportation Research Part C: Emerging Technologies}

\opening{Dear Editors,}

Please consider my manuscript, ``\textbf{ShiftGuard: Risk-Controlled
Spatiotemporal Traffic Forecasting under Compound Regime Shifts and Sensor
Outages},'' for publication as a full-length research article in
\textit{Transportation Research Part C: Emerging Technologies}.

Network-wide speed forecasts feed route guidance, ramp metering, and incident
management, yet the benchmarks used to compare forecasting models assume that
every detector reports and that traffic conditions resemble the training
period. Traffic management centers face the opposite: loop detectors fail
singly, along corridors, or across the whole network, often while traffic is
moving between free flow and congestion. The manuscript asks a question that
matters to such centers: when detectors fail and conditions shift, can a
forecaster still be accurate, report honest uncertainty, and tell the operator
when its forecast should not be used?

The manuscript makes three contributions.
\begin{enumerate}
  \item A leakage-controlled stress protocol that crosses independent,
  temporal-burst, contiguous-region, and high-centrality hub detector outages at
  five severities with traffic regimes on METR-LA and PEMS-BAY. The same saved
  outage masks are used for every model and seed, and the region and hub
  outage types are never seen during training.
  \item A model-agnostic reliability recipe: training on outage-shaped
  observation masks, a Laplace median--scale output, rolling node--horizon
  conformal calibration, and a forecast-level act/abstain rule whose threshold
  is chosen and certified on disjoint past data.
  \item Evidence that the recipe, not the architecture, carries the
  robustness. Standard baselines from DCRNN to STAEformer are retrained with the
  identical recipe and all become more robust; ablations of the compact
  reference model, ShiftGuard, are reported in full, including the components
  that did not help.
\end{enumerate}

Under synthetic outages averaged over four outage types and five severities,
ShiftGuard lowers mean absolute error relative to the strongest locally
reproduced transformer, STAEformer, from 3.74 to 3.27 mph on METR-LA and from
2.49 to 1.87 mph on PEMS-BAY. At 30\% severity it releases 91\% and 98\% of
forecasts with 7\% and 2\% failure, while using about 14 times fewer parameters
than STAEformer. Most of that gap comes from the recipe: retrained with the
same outage augmentation, every baseline's error under outages falls by
\provisional{10--29\%}, and STAEformer becomes slightly more accurate than
ShiftGuard. I report this openly because it is the point of the paper: the
recipe transfers to existing forecasters, and ShiftGuard offers near-best
accuracy at a fourteenth of the parameters and \provisional{5--9} times less
training time. Because synthetic
outages cannot establish field behavior, the manuscript also evaluates METR-LA's
own missing detector data, without any synthetic mask. There, average accuracy
is close to STAEformer's (3.00 versus 3.03 mph), but error on detectors that
were silent for the whole input window is 10--22\% lower, 90\% prediction
intervals keep their coverage in every missingness group, and 98.8\% of
forecasts are released with 5.1\% failure, compared with 79.8\% and 6.2\% for
STAEformer. The analysis also reports a failure mode openly: during
network-wide detector dropouts, every model, including ShiftGuard, exceeds the
10\% risk target, which motivates an explicit fallback in deployment.

The claims are deliberately bounded. ShiftGuard is not a clean-data state of the
art, the reference model was selected on validation data only, and the
manuscript does not claim operational dependability established by field
deployment. I believe this candor, together with the reusable protocol and
recipe, suits the journal's interest in emerging technologies whose
reliability can be assessed before deployment. The work should interest readers
working on traffic forecasting, detector data quality, uncertainty
quantification, and decision support in traffic management.

The implementation, configurations, fixed seeds, outage generators, and
table-generation scripts are available at
\href{https://github.com/shishirbit/ShiftGuard-ITS}{github.com/shishirbit/ShiftGuard-ITS},
and the manuscript gives official access information for both datasets.

I confirm that this manuscript is original, has not been published previously,
and is not under consideration by another journal. I am the sole author and
have approved its submission. The study uses previously published benchmark
data and involves no human participants, human samples, or animals. I have no
competing interests to declare.

Thank you for considering this work.

\closing{Sincerely,}

\end{letter}
\end{document}
